\documentclass[12pt]{report}

\usepackage{graphicx}
\usepackage[spanish,mexico]{babel}
\usepackage[utf8]{inputenc}
\usepackage{amsmath}
\usepackage{fancyhdr}

\graphicspath{ {./images/} }

\setlength{\oddsidemargin}{0in}
\setlength{\textwidth}{6in}
\setlength{\topmargin}{0in}
\setlength{\voffset}{-0.5in}
\setlength{\hoffset}{0.5in}
\setlength{\textheight}{9in}
\setlength{\textwidth}{6in}
\setlength{\topskip}{0in}
\setlength{\parskip}{2ex}

\title{
   {“Diseño y desarrollo de sistema multi-plataforma para gestión de ventas y logística de microempresa.”}\\
   {\large INSTITUTO TECNOLÓGICO DE TIJUANA}\\
   {\includegraphics{escudo1.png}}
}
\author{José Luis Murillo Ríos}
\date{21 DE JUNIO DEL 2019}


\begin{document}
\setcounter{page}{1}
\pagenumbering{roman}
\thispagestyle{empty}
\begin{center}
   INSTITUTO TECNOLÓGICO DE TIJUANA\\[0.75cm]
   Ingeniería en Sistemas Computacionales\\[0.8in]
\end{center}
\begin{figure}[h]
\begin{center}

\includegraphics[height=5.5 cm]{escudo1.png}
\vspace{0cm}
\end{center}
\end{figure}
\vspace{1cm}
\begin{center}
DISEÑO Y DESARROLLO DE SISTEMA MULTI-PLATAFORMA PARA GESTIÓN DE VENTAS Y LOGÍSTICA DE MICROEMPRESA.\\[4mm]
POR\\[4mm]
JOSÉ LUIS MURILLO RÍOS\\[1cm]
\small PARA OBTENER EL TÍTULO DE INGENIERO EN SISTEMAS COMPUTACIONALES
\vfill
TIJUANA, B.C. \hfill AGOSTO 2019
\end{center}

\newpage
\thispagestyle{empty}
\begin{center}
DISEÑO Y DESARROLLO DE SISTEMA MULTI-PLATAFORMA PARA GESTIÓN DE VENTAS Y LOGÍSTICA DE MICROEMPRESA.\\[1.3cm]
POR\\[0.3cm]
JOSÉ LUIS MURILLO RÍOS\\[0.3cm]
\small AGOSTO 2019\\[0.7cm]
\end{center}

% \chapter*{RESUMEN}
\newpage

\begin{center}
  {\Large \bf{RESUMEN}}
\end{center}
Cada día mas emprendedores inician su negocio con la ayuda de servcicios de venta en linea; muchas de estas empresas 
tienden a crecer y requerir de sistemas mas especializados que se ajusten a necesidades particulares.\\[0.8cm]
Se desarrolla un sistema que tiene como objetivo el llevar control de los procesos de venta y logística de una empresa.
El propósito del proyecto es integrar los servicios existentes de venta en linea de la compañía
y extender su uso, adaptándalos a los requerimientos específicas del cliente, así como diseñar una interfáz gráfica que permita
proporcionar un servicio mas rápido y agilice los proceso de elaboracion y entrega de productos.
 
% \chapter*{DEDICATORIA}
\newpage
\begin{center}
  {\Large  \bf{DEDICATORIA}}
\end{center}
\begin{center}
  A mis abuelos, por creer siempre en mi\\
  Gracias
\end{center}

 
% \chapter*{AGRADECIMIENTOS}
\newpage
\begin{center}
  {\Large  \bf{AGRADECIMIENTOS}}
\end{center}
\begin{center}
  Quiero dar gracias a Bolt Media por incluirme en su equipo y darme la oportunidad de crecer.
\end{center}


\newpage
\tableofcontents

\newpage
\listoftables
\addcontentsline{toc}{chapter}{Índice de Tablas} %%% Para introducir una línea en el índice
\listoffigures
\addcontentsline{toc}{chapter}{Índice de Figuras}

\chapter{Introducción}
\thispagestyle{fancy}
\pagenumbering{arabic} %%% esto es para regresar el modo de numeración a numeración arábiga
\setcounter{page}{1}  %%% empezamos en página 1
\thispagestyle{empty}  %%%% la primera página no va enumerada

\section{Antecedentes y definición del problema}
Se solicitó a la empresa Bolt Media Internacional S. de R.L. de C.V. desarrollar una plataforma que conectara las transacciones de rosesland.com, una plataforma de ventas en linea desarrollada en Shopify a las operaciones internas de la florería, como lo son las ventas de mostrador, la manufactura y la entrega de los productos.
\vspace{0.8cm}
Actualmente la florería realiza todo este proceso con comandas que los empleados de ventas llenan a mano, para después enviarlos al área de elaboración de arreglos florales y posteriormente se entregan al encargado de logística para dar indicaciones de la distribución de los productos. Este proceso de venta es lento y genera una experiencia poco placentera para los compradores, a su vez el área de manufactura está conformado por artesanos y trabajadores del campo que tienen un nivel de comprensión de lectura bajo y pueden llegar a cometer errores que retrasen todos los procesos. Por tales motivos este sistema debe ser fácil de usar para usuarios de distintas áreas y a su vez el diseño debe de adaptarse tanto a diferentes tamaños de pantallas como a diferentes dispositivos móviles. Se requiere una base de datos que notifique en tiempo real los cambios
en la información asi como un servidor que administre los permisos adecuados para su manipulación.

\section{Motivación para atenderlos}
Implementar un sistema diseñado principalmente para ser utilizado en pantallas táctiles
que muestre las transacciones realizadas durante el día tanto en mostrador como en la pagina web,
dando un informe detallado de las operaciones de la compañia.

\section{Objetivos generales}
Automatizar el proceso de ventas y entregas de la floreía Rosesland 
e integrar los servicios de su tienda en linea rosesland.com 
mediante una aplicación web progresiva, que sea compatible con los dispositivos existentes de la empresa
y que funcione adecuadamente en navegadores de internet modernos.

\section{Objetivos específicos}
\begin{enumerate}
\item Desarrollar un servidor que adminstre la autenticación de los empleados y su acceso a la información.
\item Implementar una base de datos que notifique a los usuarios cambios en la información en tiempo real.
\item Diseñar una aplicacion web progesiva que garantice una experiencia de usuario óptima para todas las areas de la empresa
\item Integrar los servicios de la tienda en linea rosesland.com con las operaciones internas de la empresa, uniendo los procesos en una sola plataforma
\end{enumerate}

\section{E-commerce (Comercio electrónico)}
El comercio electrónico se refiere al proceso de compra o venta de productos o servicios a través de Internet. 
Las compras en línea se están volviendo cada vez más populares debido a la velocidad y facilidad de uso para los clientes. 
Las actividades de comercio electrónico, como la venta en línea, pueden dirigirse a consumidores u otras empresas. 
% Business to Consumer (B2C) implica la venta en línea de bienes, servicios y el suministro de información directamente a los consumidores. 
% Business to Business (B2B) se refiere al intercambio en línea de productos, servicios o información entre empresas.
Vender en línea puede ayudar a su empresa a llegar a nuevos mercados y aumentar sus ventas e ingresos 
(ya sea a través de su propio sitio web o de un sitio de mercado electrónico).
\subsubsection{Shopify}
Shopify es un servicio web que le permite configurar una tienda en línea para vender sus productos. 
Le da la facilidad organizar sus productos, personalizar el diseño de su tienda, 
aceptar pagos con tarjeta de crédito, rastrear y responder a pedidos. 
Shopify.com permite a los vendedores elegir entre opciones de diseño gratuitas 
o diseños personalizados creados por los usuarios.

\section{Base de datos NoSQL}
Las bases de datos NOSQL son una alternativa emergente a las bases de datos relacionales más utilizadas. 
Como su nombre lo indica, no reemplaza completamente a SQL, 
sino que lo complementa de tal manera que puedan coexistir. \\[0.3cm]
El concepto de NOSQL se desarrolló hace mucho tiempo, 
pero fue después de la introducción de la base de datos como servicio (DBaaS) 
que obtuvo un reconocimiento destacado. Debido a la alta escalabilidad proporcionada por NOSQL, 
fue visto como un importante competidor del modelo de base de datos relacional. 
A diferencia de RDBMS, las bases de datos NOSQL están diseñadas para escalar fácilmente a medida que crecen. 
La mayoría de los sistemas NOSQL han eliminado el soporte multiplataforma 
y algunas características adicionales innecesarias de RDBMS, 
haciéndolos mucho más livianos y eficientes que sus contrapartes RDMS.
\subsubsection{Cloud Firestore}
Cloud Firestore es el servicio de base de datos de Google Firebase para aplicaciones móviles. 
Cloud Firestore almacena datos en documentos organizados en colecciones. 
Los datos simples se almacenan en documentos, 
lo cual es fácil y similar a la forma en que se almacenan los datos en JSON. 
Los datos jerárquicos complejos se organizan convenientemente a escala utilizando subcolecciones 
dentro de los documentos. \\[0.3cm]
La escalabilidad es completamente automática, 
lo que significa que no es necesario compartir sus datos en varias instancias. 
Los cargos de Cloud Firestore se basan en las operaciones realizadas en su base de datos 
(lectura, escritura, borrado), ancho de banda y almacenamiento. 
Admite límites de gasto diario para proyectos de Google App Engine, 
para garantizar que no exceda los costos con los que el usuario se sienta cómodo.

\end{document}